\begin{afterword}
\summaryhead

The author reports that when the subject of advanced AI comes up, people often answer with \textsc{The Terminator}, \textsc{The Matrix} or Asimov's robots. The post argues that this is an error even though no one believes the films are true.

Science fiction is not forecasting. Stories are shaped by what makes a good story, and a story must follow one line of events, full of details, with no way to say ``I don't know.'' Starting from a film skips the hard work of finding which possibilities deserve attention; it frames the debate as the film frames it, as a struggle of us against them; and it leaves out what matters about the subject. Brief primes affect judgment, so a two-hour film should affect it more. People do not believe films, but recall them as if they were historical cases from somewhere else. The author gives an example from a conversation in which the author argued from a character in a Vernor Vinge novel, and then noticed. Seeing a few Borg creates a category, as seeing three zebras does. Worst of all, in the author's estimation, other people's imaginings stop us from using our own, as Orwell said of stock phrases and Pirsig of a blocked student.

\responsehead

The post's warning is sound. A film is built to be watched, not to be right, and a vivid story can set the terms in which a question is asked. The parody of a story told in probabilities makes the point well, and the Vinge episode, in which the author catches the error in the author's own words, is the post's best evidence.

The rest of the support is thinner than it could have been. How often the error occurs rests on the author's conversations (``more than half the time,'' ``So far as I can tell''). The one experimental claim, that a glimpse of ``a single word'' strongly shifts probability estimates, has no source; the author's own chapter on global risks rested the argument on anchoring experiments instead. The ancestral-environment sentence repeats ``Making History Available'' without evidence. Bostrom's good-story bias, which Bostrom offered as a supposition, is introduced as something Bostrom ``points out.'' Meanwhile there was published research, such as Green and Brock's (2000), that supported the post's central claim, that fiction shapes belief even when it is known to be fiction, and the post does not use it.

The framing argument names its own alternative. The film, it says, brings to mind robot armies instead of ``a superintelligence snapping nanotechnological fingers.'' That alternative is offered only as what the frame leaves out, but it is also a specific and vivid picture, and the post gives it no support.

Finally, the claim the post calls the most damaging is offered as an estimate, supported by two quotations about something else, stock phrases in Orwell and received opinion in Pirsig, and then restated in the last line without the hedge.

In short: a sound warning about letting stories frame a forecast, supported by the author's conversations instead of the research available.
\end{afterword}
